

%% ===================================================================
\section{\vrev{한계점}}
\label{sec:limitations}
%% ===================================================================

본 연구는 이론적 모델 수립과 시뮬레이션 기반 검증에 초점을 맞추고 있으며,
다음의 한계를 지닌다. \p{[}표~\vref{tab:ch6-limitations}\p{]}은 각 한계와
대응하는 향후 연구 방향을 종합한다.

\begin{table}[htbp]
\centering
\caption{\revdel{한계점 및 향후 연구 방향}\,\revadd{한계점 종합}}
\label{tab:ch6-limitations}
\small
\begin{tabularx}{\textwidth}{c l X X}
\toprule
\textbf{\#} & \textbf{한계 영역} & \textbf{한계 내용} & \textbf{향후 연구 방향} \\
\midrule
1 & 시뮬레이션 기반 검증 &
  $n = 10{,}000$ MC 시뮬레이션과 $10^8$회 다중 시드 검증은 통계적으로 강건하나,
  실제 조직 환경에서의 현장 실험(field test)을 수반하지 않음 &
  제조$\cdot$금융$\cdot$IT 서비스 등 산업별 파일럿 현장 검증, 실제 장애 이벤트에서의
  $\eta$ 예측 정확도 평가 \\
\addlinespace
2 & 고정 $R_{\max}/R_0$ 비율 &
  $R_0 = 6.0$h, $R_{\max} = 24.0$h ($1:4$ 비율) 고정,
  산업$\cdot$조직별 비율 변동 시 모델 행태 미탐색 &
  $R_{\max}/R_0$ 비율을 $[2, 8]$ 범위에서 변동시키는 민감도 분석,
  동적 $R_0(t)$ 모형으로의 확장 \\
\addlinespace
3 & 분포 가정 &
  C2: $U \sim N(3,1)$, C3: $6 \cdot \mathrm{Pareto}(\alpha=3)$ 고정,
  실제 불확실성은 다양한 꼬리 두께를 가질 수 있음 &
  \jrev{일반화 파레토 분포(GPD)와 GEV 분포 도입, 가우시안-파레토 혼합 분포 모형 개발, $\alpha$ 변동에 따른 민감도 분석(P85.8 교차점 이동 및 감쇠 비율 변화)} \\
\addlinespace
4 & 시간 동역학 부재 &
  현행 모델은 시점별 단면(cross-sectional) 분석,
  시간에 따른 $O(t)$, $U(t)$의 자기상관, 추세, 계절성 미반영 &
  \jrev{시계열 모형(ARIMA, GARCH) 통합, 베이지안 온라인 학습 기반 자기 교정 메커니즘 구축, 복구 과정 중 실시간 $\DRTO$ 갱신 프로토콜 수립} \\
\addlinespace
\tmp{5} & 심층 불확실성 미반영 &
  현행 모델은 분포 유형(가우시안/파레토)이 사전에 특정 가능하다는 전제,
  모델$\cdot$시나리오$\cdot$확률 모두 불확실한 심층 불확실성(deep uncertainty) 상황은
  적용 범위를 벗어남 &
  로버스트 최적화, 시나리오 발견,
  info-gap 의사결정 이론 등 분포 비의존적 접근법과의 통합 탐색 \\
\addlinespace
\tmp{6} & 국가 위기경보 체계와의 연계 부재 &
  $\eta$ 기반 에스컬레이션은 정량적 지표에 의해 등급이 자동 결정되나,
  국가 위기경보(관심--주의--경계--심각)의 질적 판단 체계와의
  공식적 연계 메커니즘이 마련되지 않음 &
  국가 위기경보 4단계와 $\eta$ 임계값의 매핑 프레임워크 개발,
  질적 판단과 정량적 지표의 상호 보완적 의사결정 모형 구축 \\
\addlinespace
\jrev{7} & \jrev{\frev{분포 환경 적합성의 간접 시사}} &
  \jrev{본 연구는 \chg{\citet{lee2026drto}}의 $\kappa^{*}=1.2$를 C2/C3 조건에 적용한 후 9개 가설 PASS의 통합 해석으로 \frev{운영 적합성을 시사하였으나}, 직접적 격자 탐색은 4기준 종합점수 $F(\kappa)$가 C1 특화 지표라 C2/C3에 부적합} &
  \jrev{환경 무관 다기준 종합점수 재정의, C2/C3 조건에 적합한 새로운 4기준 도출 후 직접 격자 탐색 수행} \\
\bottomrule
\end{tabularx}
\end{table}


%% -------------------------------------------------------------------
\subsection{\vrev{시뮬레이션 기반 검증의 한계}}
\label{subsec:lim-simulation}
%% -------------------------------------------------------------------

$n = 10{,}000$ 몬테카를로 시뮬레이션과 $10^8$회 다중 시드 검증은
통계적으로 강건한 결과를 제공하나, 실제 조직 환경에서의 현장 실험을
수반하지 않았다. 시뮬레이션에서 확인된 효과 크기
(Cohen's $d = -1.180$, $+1.871$)와 이중채널 감쇠($0.52$배)가
실제 복구 과정에서도 동일한 크기로 발현되는지는
추가적인 실증 연구를 통해 확인되어야 한다.
특히 관측성 지표 $O(t)$의 실시간 측정 정확도, 불확실성 분포의
실제 형태, 그리고 조직의 복구 역량 변동이 시뮬레이션 가정과
어떻게 괴리되는지를 현장 데이터로 검증하는 것이 필요하다.
$O_{\text{raw}}$의 조작적 정의로서, 장애 탐지율, 로그 가시성, 복구 진척 추적 수준 등의
하위 지표를 리커트 척도로 평가하여 $[0, 1]$로 정규화하는 방법을 예시할 수 있으나,
측정 지표의 가중치 산정과 \chg{산업별 및 조직별} 특성을 반영한 측정 도구의 정밀화는
후속 연구에서 수행되어야 한다.


%% -------------------------------------------------------------------
\subsection{\vrev{고정 $R_{\max}/R_0$ 비율의 한계}}
\label{subsec:lim-ratio}
%% -------------------------------------------------------------------

$R_{\max}/R_0$ 비율을 모든 산업에서 4로 고정하여 비율 변동의 영향을 탐색하지
않은 점은 본 연구의 한계이다. 본 연구에서는 $R_0 = 6.0$h, $R_{\max} = 24.0$h로 $R_{\max}/R_0 = 4$의
고정 비율을 사용하였다. 이 비율은 물류 산업의 전형적 설정(\p{[}표~\vref{tab:ch6-industry-params}\p{]})과
일치하나, 금융($R_{\max}/R_0 = 4$), IT 서비스($R_{\max}/R_0 = 4$) 등
모든 산업에서 동일 비율이 적용되었다. 실무에서는 \chg{산업과 조직, 규제} 환경에 따라
이 비율이 $2$에서 $8$ 이상까지 변동할 수 있으며, 비율 변동이 시그모이드의
작동 영역($z$ 범위)과 이중채널 감쇠 비율에 미치는 영향은
체계적으로 탐색되지 않았다.

향후 연구에서는 $R_{\max}/R_0$ 비율을 독립 변수로 설정한 민감도 분석을 통해,
비율 변동에 대한 모델의 강건성을 검증하고, 산업별 최적 비율 범위를
도출하는 것이 필요하다. 나아가 $R_0(t)$ 자체를 시간 변동 함수로 확장하여,
BIA 갱신 시마다 기준선이 적응적으로 조정되는 동적 기준선 모형을 개발할 수 있다.


%% -------------------------------------------------------------------
\subsection{\vrev{분포 가정의 한계}}
\label{subsec:lim-distribution}
%% -------------------------------------------------------------------

불확실성 분포를 가우시안과 파레토($\alpha = 3$)의 두 형태로 고정한 점은 본 연구의
또 다른 한계이다. C2 조건에서의 불확실성 분포 $U \sim N(3,1)$과 C3 조건에서의
$6 \cdot \mathrm{Pareto}(\alpha = 3)$은 각각 경꼬리와 중꼬리의
대표적 형태를 포착하기 위한 선택이다. 그러나 실제 재난 복구 과정에서의
불확실성은 $\alpha < 3$(더 무거운 중꼬리)이나 $\alpha > 3$(더 가벼운 경꼬리)의
다양한 특성을 보일 수 있다.

$\alpha$ 값의 변동에 따른 핵심 결과의 민감도\chg{(}P85.8 교차점의 이동,
CVaR 격차의 변화, 이중채널 감쇠 비율\chg{)}는 향후 연구에서 체계적으로 탐색되어야 한다.
일반화 파레토 분포(GPD)나 극단값 이론(EVT)의 GEV 분포를
불확실성 모형에 도입하여, 다양한 꼬리 두께에 대한 모델의 적응성을
확보하는 것이 바람직하다.

\jrev{또한 본 연구의 \textbf{\frev{분포 환경 적합성} 종합 명제}(4.13절)는 \chg{\citet{lee2026drto}}의
$\kappa^{*}=1.2$의 C2/C3 조건 적용성을 9개 가설(\jrev{H2--H10}) PASS의 통합 해석으로
\frev{시사한다}. 이는 직접적 격자 탐색의 대안으로서 학술적 정직을 우선한 접근이지만,
\chg{\citet{lee2026drto}}의 4기준 종합점수 $F(\kappa) = f_1 \cdot f_2 \cdot f_3 \cdot f_4$
(\chg{실무 유의성, 비포화, 효과 크기, 설명력})가 C1 조건 기반으로 정의되어
C2/C3 환경에서 직접 격자 탐색에 부적합하다는 \textbf{지표 일반화의 한계}를
내포한다. 환경 무관 다기준 종합점수의 재정의 및 C2/C3 환경에서의
직접 격자 탐색은 본 연구 범위를 벗어나는 후속 연구 과제이다.}


%% -------------------------------------------------------------------
\subsection{\vrev{시간 동역학의 부재}}
\label{subsec:lim-temporal}
%% -------------------------------------------------------------------

현행 $\DRTO$ 모델은 특정 시점의 $O(t)$와 $U(t)$를 입력으로 받아
해당 시점의 $\DRTO$를 산출하는 단면(cross-sectional) 모형이다.
시간에 따른 입력 변수의 자기상관(autocorrelation), 추세(trend),
계절성(seasonality)은 모형에 내생적으로 반영되지 않는다.

실제 복구 과정에서는 $O(t)$가 복구 진행에 따라 점진적으로 증가하고,
$U(t)$는 초기에 높았다가 정보 획득에 따라 감소하는 시간적 패턴을 보인다.
이러한 동역학을 포착하기 위해, 베이지안 온라인 학습(Bayesian online learning)을
적용하여 장애 이벤트 관측 시마다 매개변수를 적응적으로 갱신하는
자기 교정(self-calibrating) 메커니즘의 구축이 향후 연구 방향이다.
ARIMA, GARCH 등 시계열 모형~\citep{durbin2012timeseries,hyndman2018forecasting}과의 통합을 통해
$\DRTO$ 모형에 시간 차원을 추가하는 확장도 고려할 수 있다.


%% -------------------------------------------------------------------
\subsection{\vrev{심층 불확실성의 미반영}}
\label{subsec:lim-deep-uncertainty}
%% -------------------------------------------------------------------

심층 불확실성을 반영하지 못하는 점은 현행 모델의 구조적 한계이다.
현행 $\DRTO$ 모델은 불확실성의 확률분포 유형이 사전에 특정 가능하다는 전제 하에
설계되었다. C2 조건에서는 가우시안, C3 조건에서는 파레토 분포를 채택하여
각각 경꼬리와 중꼬리 특성을 포착한다. 그러나 팬데믹, 기후변화에 의한 복합재난,
전례 없는 대규모 사이버 공격 등 \textbf{심층 불확실성}(deep uncertainty) 상황에서는
분포의 형태는 물론 모델 구조와 시나리오 자체를 사전에 특정할 수 없다.

이는 \subsecref{subsec:uncertainty-risk}에서 정리한 불확실성 분류 체계의
네 번째 유형\chg{(}\chg{모델, 시나리오, 확률} 모두 불확실한 심층 불확실성\chg{)}에 해당하며,
현 프레임워크의 구조적 적용 한계를 구성한다.
나이트적(Knightian) 불확실성은 파레토 분포로 보수적 근사가 가능하나,
심층 불확실성은 분포 선택 자체가 불가능하므로 확률분포 기반 모형의 범위를 벗어난다.

향후 연구에서는 로버스트 최적화(robust optimization),
시나리오 발견(scenario discovery),
info-gap 의사결정 이론 등 분포 비의존적(distribution-free) 접근법을
$\DRTO$ 프레임워크와 통합하여,
심층 불확실성 상황에서도 복구시간의 동적 조정이 가능한 확장 모형을 탐색할 필요가 있다.


%% -------------------------------------------------------------------
\subsection{\vrev{국가 위기경보 체계와의 연계}}
\label{subsec:lim-national-alert}
%% -------------------------------------------------------------------

본 연구의 $\eta$ 기반 에스컬레이션은 국가 위기경보 체계와의 공식적 연계
메커니즘을 아직 갖추지 못하였다. \subsecref{subsec:escalation-thresholds}에서 논의한 바와 같이,
현행 국가 위기경보 체계(관심--주의--경계--심각)는
재난 및 안전관리 기본법에 근거한 \textbf{질적 판단}에 기반하여
전문가 합의와 상황 평가를 통해 등급이 결정된다.
반면 본 연구의 $\eta$ 기반 에스컬레이션은
$\DRTO / R_0$라는 \textbf{정량적 지표}에 의해 등급이 자동으로 산출되므로,
두 체계는 상호 보완적 관계를 형성할 잠재력을 지닌다.

구체적으로, 국가 위기경보가 발령된 상황에서 $\eta$ 값을 병행 산출하면
질적 판단의 객관적 근거를 제공할 수 있으며,
역으로 $\eta$의 급격한 상승이 감지될 때 위기경보 상향 검토의
조기 경고(early warning) 신호로 활용할 수 있다.
향후 연구에서는 국가 위기경보 4단계와 $\eta$ 임계값 간의
체계적 매핑 프레임워크를 개발하고,
질적 판단과 정량적 지표가 결합된 하이브리드 의사결정 모형을 구축하여
국가 수준의 재난 대응 체계에 $\DRTO$의 정량적 기여 가능성을
실증적으로 검증하는 것이 필요하다.


%% -------------------------------------------------------------------
\syncr{\subsection{하위 변수 조작화의 단순화}}
\syncr{\label{subsec:lim-operationalization}}
%% -------------------------------------------------------------------

\syncr{본 연구는 관측성($O_{\text{raw}}$, $k_O$)과 불확실성($U$)의
제1계층 하위 지표까지의 조작화 \textbf{예시}를 \subsecref{app:operationalization}에
참고용으로 제시하였으나, 가중치($w_i$, $\alpha$, $\beta$)의 산업별 튜닝,
원시 지표의 측정 도구별 \chg{정확도와 편향} 검증, 현장 데이터 기반 \chg{수렴 및 예측}
타당도 확보는 본 연구의 범위를 벗어나는 후속 연구 과제이다. 제시된
조작화 예시와 가상 기업 라이프사이클은 실무 적용의 방향성을 보여주기
위한 참고용일 뿐 검증된 방법론이 아니며, 심리측정학적 엄밀성을 갖춘
측정 도구 개발 및 종단 연구를 통한 예측 타당도 확보가 선결되어야 한다.}

\vspace{1em}

\fpara{본 연구의
가장 중요한 학술 한계는 \emph{$\kappa = 1.2$의 환경 한정성}이다. \chg{\citet{lee2026drto}}의 도출값
$\kappa = 1.2$는 C1 조건(관측성 단독, 가우시안 베이스라인) 격자 탐색에서 도출되었으며,
C2/C3 환경에서의 적합성은 4.13절의 H2--H10 PASS 통합 해석으로 \emph{운영 적합성을
시사}하나 \emph{환경별 최적성을 입증}하지는 못한다. 따라서 본 연구는 $\kappa = 1.2$를
BCM 일반 운영의 \emph{\chg{잠정적이고 보수적인} 실무 기본값}(provisional and conservative operational
default)으로 채택하며, \emph{환경 무관 보편 상수성을 주장하지 않는다}. 환경 무관 $\kappa$의
\emph{존재 여부}에 대한 다방법 탐색은 후속 연구 논문의 핵심 의제로 분리된다
(5.5.3절 향후 연구 참조).}

본 장에서 제시한 \chg{이론적, 방법론적, 실무적} 시사점과 한계는
$\DRTO$ 프레임워크의 현재 위상을 명확히 규정하는 동시에,
후속 연구의 구체적 방향을 제시한다. 특히 2단계 프레임워크(Core C2 + Tail C3),
$\eta$ 기반 에스컬레이션 체계, 6개 산업 적용 가이드라인은
$\DRTO$의 이론적 기여를 실무적 가치로 전환하기 위한 구체적 기반을 제공하며,
\tmod{\chref{ch:discussion}}{\tsecref{sec:conclusion-future-merged}}에서는 이러한 논의를 종합하여 본 연구의 결론과 향후 연구 방향을 제시한다.% 3차 심사 (2026-05-09): 자기참조 오류 수정 — \chref{ch:discussion}가 이 5장 자체를 가리킴. 5.5절을 가리키도록 \tsecref(빨간색 secref)로 수정.
